%\documentclass[final,journal,compsoc,hidelinks]{article}
\documentclass[11pt,final,hidelinks]{article}
\makeatletter\def\input@path{{sty/}}\makeatother
\usepackage{lmodern}
\usepackage[nohints]{minitoc}
\usepackage[margin=1in]{geometry}
\usepackage[english]{babel}
\usepackage{graphicx}
\usepackage[activate={true,nocompatibility},final,tracking=true,kerning=true,spacing=true,factor=1100,stretch=10,shrink=10]{microtype}
\microtypecontext{spacing=nonfrench}
\usepackage[noend]{algorithm2e}
\usepackage{caption}
\usepackage{mathtools}
\usepackage{commath}
\usepackage{minted}
\usepackage{enumerate}
\usepackage{amsmath}
\usepackage{subfiles}
\usepackage{booktabs}
% Phase B (2026-06-08): the glossaries package was loaded [xindy,toc,acronyms,symbols]
% with \printglossary but no entries were ever defined, yielding an empty glossary and an
% extra build step. Removed (package load and the \printglossary call) until entries exist.
\usepackage{latexsym}
\usepackage{amsthm}
\usepackage{amssymb}
\usepackage{pgfplots}
\usepackage{wasysym}
\usepackage{hyperref}
\usepackage{tikz}
\usepackage{tikzscale}
\usepackage[square,numbers]{natbib}
\bibliographystyle{plainnat}
\usepackage[utf8]{inputenc}
\usepackage[T1]{fontenc}
\usepackage{cleveref}
\usepackage[super]{nth}
\usepackage{siunitx}
\usepackage[section]{placeins}
\usepackage{subcaption}

%\usepackage{alexmisc}
\usepackage{functionnotation}
\usepackage[plain]{setnotation}
\usepackage{approxsetnotation}
\usepackage{probabilitynotation}
\usepackage[fancy]{relationnotation}
\usepackage{approxrelationnotation}
\usepackage[section]{envnotation}
%\usepackage{algorithmnotation}
%\usepackage{alexmisc}

\hypersetup{
    hypertexnames=false,
    pdftitle={Error contracts for approximate relations},
    pdfauthor={Alexander Richard Towell},               % author
    pdfsubject={computer science},              % subject of the document
    pdfkeywords={
        probabilistic data structure,
        partial function,
        random approximate set,
        random approximate data type,
        random approximate relation,
        random approximate map,
        relation,
        set        
	},                    % keywords
    colorlinks=false,                           % false: boxed links;
    citecolor=green,                            % color of links to 
    filecolor=magenta,                          % color of file links
    urlcolor=green                              % color of external
}

\PassOptionsToPackage{unicode}{hyperref}
\PassOptionsToPackage{naturalnames}{hyperref}
\pgfplotsset{compat=1.14} 

\graphicspath{{img/}}


\title{Error contracts for approximate relations}
\author{Alexander Richard Towell\\
Southern Illinois University Edwardsville\\
\texttt{lex@metafunctor.com}\\
\small ORCID: 0000-0001-6443-9897}
\date{October 2026\\\small Final preprint version}

\begin{document}
\maketitle
\begin{abstract}
A compact approximate relation can preserve a structural constraint while
changing the dependence between query operands. Per-tuple false-positive and
false-negative rates therefore do not suffice to specify query error.
We separate exact truth, construction randomness, and frozen source identity,
and audit the resulting operator and invariant-preservation formulas.
For a monotone query with independent one-sided source errors, the sharp
worst error over exact truth profiles is $1-(1-\alpha)^{s_0}$ for false
positives and $1-(1-\beta)^{s_1}$ for false negatives, where $s_0,s_1$ are
one-sided sensitivities. These are elementary consequences of established
certificate theory. Marginal-only bounds and a conditional shared-rate
extension state what remains sound when independence is unavailable.
For normalized binary-witness queries of source occurrence at most two,
the required sensitivities have a linear-time graph formula. A stored-filter
experiment uses these contracts to calibrate symmetric relation queries,
comparing canonical keys with ordered keys followed by symmetric repair.
We distinguish exact model guarantees from empirical hash behavior and
from existing probabilistic lineage and inference methods. The resulting
framework supports a narrow, reproducible query-calibration study; it does
not establish a universal probability calculus or cryptographic security.
\end{abstract}

% Status and errata box (October 2026 freeze). \input right after \end{abstract}.
\begingroup
\setlength{\fboxsep}{7pt}%
\begin{center}
\fbox{\begin{minipage}{0.93\linewidth}\small
\setlength{\parskip}{4pt}%
\setlength{\parindent}{0pt}%
\textbf{Status and errata (October 2026).}
This is the final version of this preprint; no further revisions are planned.
It is superseded by Chapter~12 (\emph{Approximate Relations}) of the monograph
\emph{Bernoulli Types: Two-Sided Error in Approximate Membership, Retrieval,
and Computation} (A.~Towell, 2026),
\url{https://metafunctor.com/latex/bernoulli-monograph/book.html}
(project hub: \url{https://metafunctor.com/research/bernoulli/}).
Please cite the monograph in preference to this preprint.
The previous Zenodo version (v5) was titled \emph{The abstract data type of
the approximate relation}; this version is retitled to match its narrower scope.

\textbf{Corrections in this version} (relative to v5).
\begin{itemize}\setlength{\itemsep}{1pt}\setlength{\parskip}{0pt}
\item Projection: the claim that projection preserves element-wise independence
is withdrawn; projected tuples whose fibers share a source are correlated
(\cref{thm:project_sources}), and profiles must record true-member count as well as fan-in.
\item Join: dependence through a shared input is stated with its exact conditions
rather than as a universal understatement of joint error.
\item Transitivity: the claim that its preservation probability vanishes
(``fastest'') is replaced by a sufficient bound for fixed interior rates; the
empty relation with zero false-positive rate is a counterexample to unconditional vanishing.
\item The ``closure dichotomy'' is withdrawn; the retained results concern specific
orbit and simultaneous-addition constructions. The symmetric diagonal (one trial,
not two) and the use of actual rather than maximal orbit size are corrected.
\item The single-valuedness preservation probability is replaced by the exact
defined/undefined-row formula; deletion repair semantics are specified.
\item The $\log_2 m$-bit surcharge is restated as conditional rate accounting,
not a space lower bound or a price for recall.
\item Partial functions: the definition is repaired (single-valued realizations,
wrong-value errors separated from domain errors) and the unsupported
map-composition closure claim is removed.
\item The encrypted-search section no longer claims that error rates characterize
security. New in this version: query error contracts (\cref{sec:query_contracts})
and a stored-filter calibration study (\cref{sec:calibration}).
\end{itemize}

\textbf{Predecessors credited.} Probabilistic relational algebra
(Fuhr and R{\"o}lleke~\cite{fuhr1997pra}), tuple-independent query evaluation,
lineage, and provenance \cite{dalvi2007efficient,benjelloun2008uldb,green2007provenance,suciu2011},
dissociation bounds \cite{gatterbauer2014bounds}, monotone certificate and
sensitivity theory \cite{blanc2022counterfactuals}, and Bloom-filter semijoins
\cite{mullin1990semijoins}. The operator formulas and the one-sided contracts are
elementary consequences of these; no novelty is claimed for them.
A dated list of changes is in \texttt{ERRATA.md} in the source repository.
\end{minipage}}
\end{center}
\endgroup



\microtypesetup{protrusion=false}
\tableofcontents
\microtypesetup{protrusion=true}


% Fiber (pre-image) of a projected tuple under projection; see def:fiber in
% sections/rel_algebra.tex. Local to this paper.
\newcommand{\fiber}{\phi}


%\listoffigures
%\listofalgorithms


\subfile{sections/intro}
\subfile{sections/relations}
\subfile{sections/approx_rel_model}
\subfile{sections/rel_algebra}
\subfile{sections/properties}
\subfile{sections/query_contracts}
\subfile{sections/calibration}
\subfile{sections/oblivious_rel_model}
\subfile{sections/conclusion}
\bibliography{references}
\end{document}
